%
\vspace{-0.1cm}
\section{Related Work}
\label{sec:related}

The General Reinforcement Learning Architecture (Gorila) of~\cite{nair2015gorila} performs asynchronous training of reinforcement learning agents in a distributed setting.
In Gorila, each process contains an actor that acts in its own copy of the environment, a separate replay memory, and a learner that samples data from the replay memory and computes gradients of the DQN loss~\citep{mnih-dqn-2015} with respect to the policy parameters.
The gradients are asynchronously sent to a central parameter server which updates a central copy of the model.
The updated policy parameters are sent to the actor-learners at fixed intervals.
By using 100 separate actor-learner processes and 30 parameter server instances, a total of 130 machines, Gorila was able to significantly outperform DQN over 49 Atari games.
On many games Gorila reached the score achieved by DQN over 20 times faster than DQN.
We also note that a similar way of parallelizing DQN was proposed by~\cite{chavez2015deepq}.

In earlier work, \cite{li:mapreduce} applied the Map Reduce framework to parallelizing batch reinforcement learning methods with linear function approximation.
Parallelism was used to speed up large matrix operations but not to parallelize the collection of experience or stabilize learning.
\cite{grounds:parallelRL} proposed a parallel version of the Sarsa algorithm that uses multiple separate actor-learners to accelerate training.
Each actor-learner learns separately and periodically sends updates to weights that have changed significantly to the other learners using peer-to-peer communication.

\cite{tsitsiklis1994asynchronous} studied convergence properties of Q-learning in the asynchronous optimization setting.
These results show that Q-learning is still guaranteed to converge when some of the information is outdated as long as outdated information is always eventually discarded and several other technical assumptions are satisfied.
Even earlier, \cite{bertsekas1982distributed} studied the related problem of distributed dynamic programming.

Another related area of work is in evolutionary methods, which are often straightforward to parallelize by distributing fitness evaluations over multiple machines or threads~\citep{tomassini1999parallel}.
Such parallel evolutionary approaches have recently been applied to some visual reinforcement learning tasks.
In one example, \cite{koutnik2014} evolved convolutional neural network controllers for the TORCS driving simulator by performing fitness evaluations on 8 CPU cores in parallel.
